Les ions permanganates, violets, réagissent sur les ions fer(II) en milieu acide
pour les transformer en ions fer(III). L'équation associée à la réaction est~:
\begin{center}
    \ch{MnO_{4}^{-}_{(aq)} + 5 Fe^{2+}_{(aq)} + 8 H^{+}_{(aq)} ->
    Mn^{2+}_{(aq)} + 5 Fe^{3+}_{(aq)} + 4 H2O_{($\ell$)}}
\end{center}
Aux concentrations utilisées, seuls les ions permanganates sont notablement
colorés.

Dans un bécher, on introduit $V_{1}=\qty{10.0}{\mL}$ de solution de sulfate de
fer(II) de concentration $c_{1}=\qty{0.055}{\mol\per\L}$ et $V=\qty{5.0}{\mL}$
d'acide sulfurique de concentration $c=\qty{0.50}{\mol\per\L}$. On ajoute
$V_{2}=\qty{4.0}{\mL}$ de solution de permanganate de potassium de concentration
$c_{2}=\qty{0.025}{\mol\per\L}$. Le mélange est incolore.

\begin{enumerate}
    \item Faire le bilan des espèces présentes à l'état initial. Quel est le
          réactif limitant~?
    \item Construire le tableau d'avancement de la réaction.
\end{enumerate}

